% Mathématiques 5e — cours V3
% Calcul littéral et premières équations

\section{Objectifs}

\begin{itemize}
\item Produire des formules simples.
\item Calculer la valeur d'une expression littérale par substitution.
\item Tester si une égalité comportant une variable est vraie ou fausse.
\item Reconnaître si une expression est une somme ou un produit.
\item Développer et factoriser avec la distributivité simple.
\item Réduire une expression de la forme $ax+b$.
\item Démontrer une propriété générale par calcul littéral.
\item Utiliser un contre-exemple pour réfuter une assertion.
\item Formuler des conjectures avec un tableur ou un langage algorithmique.
\item Donner à la lettre le statut d'inconnue.
\item Modéliser puis résoudre des équations de type $ax=c$ ou $x+b=c$.
\end{itemize}

\section{1. Pourquoi utiliser des lettres ?}

Une lettre permet de représenter un nombre que l'on ne souhaite pas fixer immédiatement.

Elle peut représenter :
\begin{itemize}
\item une grandeur qui varie ;
\item une longueur quelconque ;
\item un nombre choisi ;
\item une inconnue que l'on cherche.
\end{itemize}

\begin{exemple}
Le périmètre d'un carré de côté $c$ s'écrit :
\[
P=4c.
\]

La formule est valable quelle que soit la valeur de $c$.
\end{exemple}

Le calcul littéral sert donc à généraliser.

\section{2. Produire une formule}

Une formule traduit une relation entre plusieurs grandeurs.

Pour un rectangle de longueur $L$ et largeur $\ell$ :
\[
P=2L+2\ell.
\]

Pour son aire :
\[
A=L\ell.
\]

Pour le successeur d'un nombre $n$ :
\[
n+1.
\]

Pour son prédécesseur :
\[
n-1.
\]

Pour son double :
\[
2n.
\]

Pour son carré :
\[
n^2.
\]

\section{3. Écriture simplifiée d'un produit}

En calcul littéral, on écrit :
\[
3\times x=3x.
\]

De même :
\[
a\times b=ab.
\]

Et :
\[
4\times(x+2)=4(x+2).
\]

\begin{attention}
On supprime le signe de multiplication uniquement lorsque l'écriture reste compréhensible.
\end{attention}

On ne transforme pas :
\[
3\times4
\]
en :
\[
34.
\]

\section{4. Calculer par substitution}

Substituer consiste à remplacer une lettre par une valeur donnée.

\begin{exemple}
Soit :
\[
E=3x+5.
\]

Pour :
\[
x=4,
\]
on écrit :
\[
E=3\times4+5.
\]

Puis :
\[
E=12+5=17.
\]

Donc :
\[
\boxed{E(4)=17}
\]
si l'on souhaite utiliser cette notation.
\end{exemple}

\section{5. Parenthèses lors d'une substitution}

Si la valeur remplaçant la lettre comporte un signe ou plusieurs termes, les parenthèses évitent les erreurs.

Pour :
\[
F=2x+1
\]
et :
\[
x=3{,}5,
\]
on écrit :
\[
F=2\times3{,}5+1.
\]

Pour une valeur comme :
\[
x=-2,
\]
on écrirait :
\[
F=2\times(-2)+1.
\]

\section{6. Tester une égalité}

Une égalité comportant une lettre peut être vraie pour certaines valeurs et fausse pour d'autres.

Considérons :
\[
2x+3=11.
\]

Pour :
\[
x=4,
\]
le membre de gauche vaut :
\[
2\times4+3=11.
\]

L'égalité est donc vraie pour $x=4$.

Pour :
\[
x=5,
\]
le membre de gauche vaut :
\[
2\times5+3=13.
\]

L'égalité est fausse pour $x=5$.

\section{7. Somme ou produit}

La nature d'une expression dépend de son opération principale.

L'expression :
\[
3x+2
\]
est une somme.

L'expression :
\[
5(x+4)
\]
est un produit.

Dans :
\[
3x+2,
\]
les termes sont :
\[
3x
\]
et :
\[
2.
\]

Dans :
\[
5(x+4),
\]
les facteurs sont :
\[
5
\]
et :
\[
x+4.
\]

\section{8. Distributivité}

\coursefigure{/images/mathematiques/5eme/calcul-litteral-distributivite.svg}{Rectangle découpé en deux parties illustrant la distributivité k fois a plus b.}{Le même rectangle peut être vu comme une aire globale k(a+b) ou comme la somme des aires ka et kb.}

\begin{propriete}
Pour tous nombres $k$, $a$ et $b$ :
\[
\boxed{k(a+b)=ka+kb}.
\]

Et :
\[
\boxed{k(a-b)=ka-kb}.
\]
\end{propriete}

Cette égalité fonctionne dans les deux sens.

\section{9. Développer}

Développer consiste à transformer un produit en somme ou différence.

Par exemple :
\[
4(x+3).
\]

On distribue le facteur $4$ :
\[
4\times x+4\times3.
\]

Donc :
\[
\boxed{4(x+3)=4x+12}.
\]

Autre exemple :
\[
5(2x-1)
=
10x-5.
\]

\section{10. Factoriser}

Factoriser consiste à transformer une somme en produit.

Par exemple :
\[
7x+21.
\]

Le facteur commun est :
\[
7.
\]

On écrit :
\[
7x+21
=
7x+7\times3.
\]

Donc :
\[
\boxed{7x+21=7(x+3)}.
\]

\section{11. Réduire une expression}

Réduire consiste à regrouper les termes de même nature.

Par exemple :
\[
3x+5x+2.
\]

On additionne les coefficients de $x$ :
\[
3x+5x=8x.
\]

Donc :
\[
\boxed{3x+5x+2=8x+2}.
\]

\section{12. Expression de la forme ax+b}

Une expression :
\[
ax+b
\]
contient :
\begin{itemize}
\item un terme dépendant de $x$ : $ax$ ;
\item un terme constant : $b$.
\end{itemize}

Par exemple :
\[
4x+7.
\]

On ne peut pas réduire :
\[
4x+7
\]
en :
\[
11x.
\]

Les deux termes ne sont pas de même nature.

\section{13. Exemple développé : développer puis réduire}

Considérons :
\[
A=3(x+4)+2x.
\]

On développe :
\[
3(x+4)=3x+12.
\]

Donc :
\[
A=3x+12+2x.
\]

On réduit :
\[
3x+2x=5x.
\]

Ainsi :
\[
\boxed{A=5x+12}.
\]

\section{14. Exemple développé : factoriser}

Considérons :
\[
B=6x+18.
\]

On écrit :
\[
6x+18
=
6x+6\times3.
\]

Donc :
\[
\boxed{B=6(x+3)}.
\]

Pour vérifier, on peut développer :
\[
6(x+3)=6x+18.
\]

\section{15. Démontrer une propriété générale}

Le calcul littéral permet de prouver une propriété valable pour tous les nombres d'un certain type.

\begin{exemple}
Montrons que la somme de deux nombres entiers consécutifs est impaire.

On note le premier entier :
\[
n.
\]

Le suivant vaut :
\[
n+1.
\]

La somme vaut :
\[
n+(n+1)=2n+1.
\]

Or :
\[
2n+1
\]
est de la forme « deux fois un entier plus un ».

C'est donc un nombre impair.
\end{exemple}

La lettre permet ici de représenter tous les cas à la fois.

\section{16. Contre-exemple}

Une seule valeur peut suffire à prouver qu'une affirmation générale est fausse.

\begin{exemple}
Affirmation : « Pour tout nombre $x$, on a $x^2>x$. »

Choisissons :
\[
x=\dfrac12.
\]

Alors :
\[
x^2=\dfrac14.
\]

Or :
\[
\dfrac14<\dfrac12.
\]

L'affirmation est donc fausse.
\end{exemple}

\begin{remarque}
Un contre-exemple ne prouve pas qu'une affirmation est vraie ; il sert à montrer qu'une affirmation générale est fausse.
\end{remarque}

\section{17. Formuler une conjecture}

Une conjecture est une affirmation que l'on pense vraie après plusieurs observations, mais qui n'est pas encore démontrée.

On peut utiliser :
\begin{itemize}
\item un tableur ;
\item un programme ;
\item plusieurs essais numériques.
\end{itemize}

\begin{attention}
Tester beaucoup de valeurs ne remplace pas une démonstration lorsqu'on affirme qu'une propriété est vraie pour tous les nombres.
\end{attention}

\section{18. La lettre comme inconnue}

Dans une équation, la lettre peut représenter un nombre inconnu.

Par exemple :
\[
x+7=15.
\]

On cherche le nombre qui rend l'égalité vraie.

Ici :
\[
x=8.
\]

En effet :
\[
8+7=15.
\]

\section{19. Équation de type x+b=c}

Pour :
\[
x+b=c,
\]
on utilise l'opération inverse de l'addition.

On obtient :
\[
x=c-b.
\]

\begin{exemple}
Résolvons :
\[
x+9=14.
\]

On calcule :
\[
x=14-9.
\]

Donc :
\[
\boxed{x=5}.
\]

Vérification :
\[
5+9=14.
\]
\end{exemple}

\section{20. Équation de type ax=c}

Pour :
\[
ax=c,
\]
avec :
\[
a\neq0,
\]
on utilise l'opération inverse de la multiplication.

On obtient :
\[
x=\dfrac ca.
\]

\coursefigure{/images/mathematiques/5eme/calcul-litteral-equation.svg}{Balance schématique représentant une égalité de type a fois x égale c.}{Résoudre une équation consiste à conserver l'égalité tout en isolant l'inconnue par une opération inverse.}

\begin{exemple}
Résolvons :
\[
4x=28.
\]

On divise par $4$ :
\[
x=28\div4.
\]

Donc :
\[
\boxed{x=7}.
\]
\end{exemple}

\section{21. Modéliser un problème par une équation}

Un abonnement coûte $5\ \text{€}$ de frais fixes et une dépense inconnue $x$.

Le total est :
\[
17\ \text{€}.
\]

On peut écrire :
\[
x+5=17.
\]

Puis :
\[
x=17-5=12.
\]

La dépense inconnue vaut :
\[
\boxed{12\ \text{€}}.
\]

\section{22. Exemple de partage}

Quatre objets identiques coûtent ensemble :
\[
26\ \text{€}.
\]

Si $x$ désigne le prix d'un objet :
\[
4x=26.
\]

Donc :
\[
x=\dfrac{26}{4}=6{,}5.
\]

Le prix d'un objet est :
\[
\boxed{6{,}50\ \text{€}}.
\]

\section{23. Vérifier une solution}

Une solution d'équation doit être vérifiée en la remplaçant dans l'égalité de départ.

Pour :
\[
3x=18,
\]
on propose :
\[
x=6.
\]

Vérification :
\[
3\times6=18.
\]

La valeur :
\[
6
\]
est bien solution.

\section{24. Programme de calcul et expression}

Programme :
\begin{enumerate}
\item choisir un nombre $x$ ;
\item le multiplier par $3$ ;
\item ajouter $4$.
\end{enumerate}

L'expression obtenue est :
\[
3x+4.
\]

Pour :
\[
x=5,
\]
la sortie vaut :
\[
3\times5+4=19.
\]

Le langage littéral permet de décrire le programme pour n'importe quelle entrée.

\section{25. Tableur et formule}

Dans un tableur, une formule peut être calculée pour de nombreuses valeurs.

On peut construire un tableau :

\begin{tabular}{c|c}
$x$ & $3x+4$ \\
$0$ & $4$ \\
$1$ & $7$ \\
$2$ & $10$ \\
$3$ & $13$ \\
\end{tabular}

On observe une régularité.

Cette observation peut conduire à une conjecture ou aider à vérifier une formule.

\section{26. Méthode de substitution}

\begin{methode}
Pour calculer la valeur d'une expression :
\begin{enumerate}
\item recopier l'expression ;
\item remplacer chaque lettre par sa valeur ;
\item utiliser des parenthèses si nécessaire ;
\item respecter les priorités opératoires ;
\item effectuer le calcul ;
\item contrôler le résultat.
\end{enumerate}
\end{methode}

\section{27. Méthode de développement}

\begin{methode}
Pour développer :
\[
k(a+b),
\]
on multiplie $k$ par chacun des termes :
\[
ka+kb.
\]

Pour :
\[
k(a-b),
\]
on obtient :
\[
ka-kb.
\]
\end{methode}

\section{28. Méthode d'équation simple}

\begin{methode}
Pour résoudre une équation arithmétique simple :
\begin{enumerate}
\item identifier l'opération effectuée sur l'inconnue ;
\item utiliser l'opération inverse ;
\item calculer la valeur de l'inconnue ;
\item vérifier dans l'égalité initiale.
\end{enumerate}
\end{methode}

\section{29. Erreurs fréquentes}

\begin{itemize}
\item Additionner des termes qui ne sont pas semblables.
\item Confondre développer et factoriser.
\item Oublier de distribuer le facteur à tous les termes.
\item Remplacer une lettre par une valeur sans parenthèses alors qu'elles sont nécessaires.
\item Croire que plusieurs essais numériques démontrent une propriété générale.
\item Confondre variable et inconnue.
\item Résoudre une équation sans vérifier la solution.
\item Transformer une égalité en modifiant un seul membre sans justification.
\end{itemize}

\section{À retenir}

\begin{aretenir}
Une lettre permet de généraliser, de modéliser ou de représenter une inconnue.

La distributivité donne :
\[
\boxed{k(a+b)=ka+kb}
\]
et :
\[
\boxed{k(a-b)=ka-kb}.
\]

Développer transforme un produit en somme.

Factoriser transforme une somme en produit.

Pour une équation :
\[
x+b=c,
\]
on utilise :
\[
\boxed{x=c-b}.
\]

Pour :
\[
ax=c,
\]
on utilise :
\[
\boxed{x=\dfrac ca}.
\]

Enfin, une propriété générale doit être démontrée ; quelques exemples ne suffisent pas.
\end{aretenir}
